\documentclass[10pt,a4paper]{article}
\usepackage[T1]{fontenc}
\usepackage[utf8]{inputenc}
\usepackage[ngerman]{babel}
\usepackage[a4paper,margin=14mm]{geometry}
\usepackage{lmodern}
\usepackage[table]{xcolor}
\usepackage{tabularx}
\usepackage{microtype}
\usepackage{amssymb}
\usepackage{tikz}
\newif\ifsolution
\solutionfalse
\pagestyle{empty}
\setlength{\parindent}{0pt}
\setlength{\parskip}{0pt}
\renewcommand{\familydefault}{\sfdefault}
\definecolor{Navy}{HTML}{17324D}
\definecolor{Blue}{HTML}{3B6EA5}
\definecolor{Gold}{HTML}{F3B33D}
\definecolor{Ink}{HTML}{1E2935}
\definecolor{Muted}{HTML}{617080}
\definecolor{Line}{HTML}{D8E1E8}
\color{Ink}
\newcommand{\kopf}[2]{%
  \colorbox{Navy}{\parbox{\dimexpr\linewidth-2\fboxsep\relax}{%
    \vspace{1.3mm}\color{white}
    \begin{tabularx}{\linewidth}{@{}Xr@{}}
      {\small\bfseries PHYSIK · KLASSE 8} &
      \colorbox{Gold}{\color{Navy}\small\bfseries\strut\ifsolution LÖSUNGEN\else #1\fi}
    \end{tabularx}
    \vspace{0.9mm}
    {\fontsize{16}{19}\selectfont\bfseries Elektrischer Strom · Widerstand und gültige Ziffern}\par
    \vspace{0.3mm}{\small #2}\vspace{1.1mm}
  }}\par\vspace{1.5mm}}
\newcommand{\frage}[2]{%
  \vspace{0.6mm}{\color{Blue}\bfseries\normalsize #1 · #2}\par
  \vspace{0.2mm}{\color{Blue}\rule{\linewidth}{0.55pt}}\par\vspace{0.4mm}}
\newcommand{\antwortfeld}[2]{%
  \vspace{0.5mm}%
  \begin{tikzpicture}[x=1mm,y=1mm]
    \path[use as bounding box] (0,0) rectangle (181.5,-7*#1);
    \ifsolution
      \node[anchor=north west,text width=179mm,inner sep=0pt,align=left,
        font={\fontsize{11}{14}\selectfont\everymath{\displaystyle}},text=Blue] at (0,-0.8) {#2};
    \else
      \foreach \i in {1,...,#1}{\draw[Line,line width=.5pt] (0,-7*\i) -- (181.5,-7*\i);}
    \fi
  \end{tikzpicture}\par}
\newcommand{\rechenfeld}[2]{%
  \vspace{0.4mm}%
  \begin{tikzpicture}[x=1mm,y=1mm]
    \path[use as bounding box] (0,0) rectangle (181.5,-#1);
    \draw[step=5mm,Line!65,line width=.25pt] (0,-#1) grid (181.5,0);
    \draw[Line,line width=.55pt] (0,0) rectangle (181.5,-#1);
    \ifsolution
      \node[anchor=north west,text width=179mm,inner sep=0pt,align=left,
        font={\fontsize{11}{14}\selectfont\everymath{\displaystyle}},text=Blue] at (0,-0.7) {#2};
    \fi
  \end{tikzpicture}\par}
\newcommand{\wahl}[2]{%
  \noindent\ifsolution\ifnum#1=1 $\boxtimes$\else $\square$\fi\else $\square$\fi\enspace #2\par}
\newcommand{\fuss}[2]{%
  \vfill{\color{Line}\rule{\linewidth}{0.5pt}}\par\vspace{1mm}
  \begin{tabularx}{\linewidth}{@{}Xr@{}}
    {\small\color{Muted}Klasse 8 · Physik · Elektrischer Strom} &
    {\small\color{Muted}\ifsolution Lösungen\else #2\fi{} · Seite #1 von 2}
  \end{tabularx}}
\begin{document}
\kopf{TEST}{Test zur Abfrage · Taschenrechner · Widerstand und gültige Ziffern}
\small\textbf{Hilfsmittel: Taschenrechner.}\par\normalsize
\frage{A1}{Widerstand, Spannung und Stromstärke berechnen}
\textbf{Berechne} in a) den Widerstand, in b) die Spannung und in c) die Stromstärke. \textbf{Notiere} jeweils Ansatz, nötige Umstellung, Einsetzen mit Einheiten, Taschenrechnerwert und passend gerundetes Ergebnis.\par
\vspace{0.5mm}
\textbf{a)} $U=7{,}35\,\mathrm{V}$; $I=0{,}42\,\mathrm{A}$\par
\rechenfeld{30}{$R=\frac{U}{I}=\frac{7{,}35\,\mathrm{V}}{0{,}42\,\mathrm{A}}=17{,}5\,\Omega$.\par $\Rightarrow R=18\,\Omega$ (zwei gültige Ziffern).}
\textbf{b)} $R=7{,}4\,\Omega$; $I=0{,}38\,\mathrm{A}$\par
\rechenfeld{30}{$R=\frac{U}{I}\Rightarrow U=R\cdot I$.\par $U=7{,}4\,\Omega\cdot0{,}38\,\mathrm{A}=2{,}812\,\mathrm{V}$.\par $\Rightarrow U=2{,}8\,\mathrm{V}$.}
\textbf{c)} $U=9{,}40\,\mathrm{V}$; $R=3{,}7\,\Omega$\par
\rechenfeld{30}{$R=\frac{U}{I}\Rightarrow I=\frac{U}{R}$.\par $I=\frac{9{,}40\,\mathrm{V}}{3{,}7\,\Omega}=2{,}540540\ldots\,\mathrm{A}$.\par $\Rightarrow I=2{,}5\,\mathrm{A}$.}
\frage{A2}{Gültige Ziffern erkennen}
\textbf{Kreuze alle richtigen Aussagen an.}\par
\wahl{1}{0,0042 A hat zwei gültige Ziffern.}
\wahl{0}{Die Null vor der 7 in 0,070 A ist eine gültige Ziffer.}
\wahl{1}{5,00 V hat drei gültige Ziffern.}
\wahl{0}{Zusätzliche Stellen im Taschenrechnerwert machen das Ergebnis genauer als die gegebenen Messwerte.}
\frage{A3}{Widerstand und Genauigkeit}
$U=8{,}25\,\mathrm{V}$; $I=0{,}46\,\mathrm{A}$; $R=\frac{U}{I}$\par
\textbf{Kreuze alle richtigen Aussagen zur Rechnung an.}\par
\wahl{0}{$17{,}9\,\Omega$ ist das passend gerundete Endergebnis.}
\wahl{1}{Bei Multiplikation und Division bestimmt die kleinste Anzahl gültiger Ziffern der gegebenen Messwerte die Anzahl im Endergebnis.}
\wahl{1}{Zwischenwerte werden nicht vorzeitig gerundet; gerundet wird erst das Endergebnis.}
\wahl{1}{Aus 8,25 V und 0,46 A ergibt sich passend gerundet $R=18\,\Omega$.}
\fuss{1}{Test}
\newpage
\kopf{TEST}{Test zur Abfrage · Taschenrechner · Widerstand und gültige Ziffern}
\frage{A4}{Größen und Einheiten}
\textbf{Ergänze} in der Tabelle jeweils die fehlende Größe, das Formelzeichen und die Einheit.\par
\vspace{0.5mm}
\renewcommand{\arraystretch}{1.45}
\begin{tabularx}{\linewidth}{|>{\raggedright\arraybackslash}X|>{\centering\arraybackslash}p{32mm}|>{\centering\arraybackslash}p{36mm}|}\hline
\rowcolor{Navy}\color{white}\bfseries Größe & \color{white}\bfseries Formelzeichen & \color{white}\bfseries Einheit \\ \hline
elektrischer Widerstand & \ifsolution $R$\fi & \ifsolution Ohm ($\Omega$)\fi \\ \hline
\ifsolution elektrische Spannung\fi & $U$ & \ifsolution Volt (V)\fi \\ \hline
\ifsolution elektrische Stromstärke\fi & \ifsolution $I$\fi & \ifsolution Ampere (A)\else A\fi \\ \hline
\end{tabularx}
\par\vspace{4mm}
\frage{A5}{Einheiten umrechnen}
\textbf{Rechne} die vier Angaben in die jeweils genannte Einheit um; erhalte dabei die Anzahl gültiger Ziffern.\par
\begin{tabularx}{\linewidth}{@{}X X@{}}
a) $125\,\mathrm{mV}=\ifsolution 0{,}125\,\mathrm{V}\else\rule{29mm}{.4pt}\,\mathrm{V}\fi$ & b) $0{,}347\,\mathrm{k\Omega}=\ifsolution 347\,\mathrm{\Omega}\else\rule{28mm}{.4pt}\,\mathrm{\Omega}\fi$ \\
\vspace{6mm}&\\[-2mm]
c) $2{,}365\,\mathrm{V}=\ifsolution 2365\,\mathrm{mV}\else\rule{28mm}{.4pt}\,\mathrm{mV}\fi$ & d) $0{,}047\,\mathrm{mA}=\ifsolution 0{,}000047\,\mathrm{A}\else\rule{28mm}{.4pt}\,\mathrm{A}\fi$
\end{tabularx}
\par\vspace{4mm}
\frage{A6}{Für die Schnellen}
\textbf{Berechne} in a) den Widerstand, in b) die Spannung und in c) die Stromstärke. \textbf{Notiere} jeweils Ansatz, nötige Umstellung, Einsetzen mit Einheiten, Taschenrechnerwert und passend gerundetes Ergebnis.\par
\vspace{0.5mm}
\textbf{a)} $U=4{,}82\,\mathrm{V}$; $I=0{,}37\,\mathrm{A}$\par
\rechenfeld{30}{$R=\frac{U}{I}=\frac{4{,}82\,\mathrm{V}}{0{,}37\,\mathrm{A}}=13{,}027027\ldots\,\Omega$.\par $\Rightarrow R=13\,\Omega$ (zwei gültige Ziffern).}
\textbf{b)} $R=18{,}6\,\Omega$; $I=0{,}27\,\mathrm{A}$\par
\rechenfeld{30}{$R=\frac{U}{I}\Rightarrow U=R\cdot I$.\par $U=18{,}6\,\Omega\cdot0{,}27\,\mathrm{A}=5{,}022\,\mathrm{V}$.\par $\Rightarrow U=5{,}0\,\mathrm{V}$.}
\textbf{c)} $U=2{,}85\,\mathrm{V}$; $R=8{,}6\,\Omega$\par
\rechenfeld{30}{$R=\frac{U}{I}\Rightarrow I=\frac{U}{R}$.\par $I=\frac{2{,}85\,\mathrm{V}}{8{,}6\,\Omega}=0{,}331395\ldots\,\mathrm{A}$.\par $\Rightarrow I=0{,}33\,\mathrm{A}$.}
\fuss{2}{Test}
\end{document}
